\documentclass[10pt,twocolumn]{article}

\usepackage[utf8]{inputenc}
\usepackage[T1]{fontenc}
\usepackage{amsmath,amssymb}
\usepackage{graphicx}
\usepackage{booktabs}
\usepackage{hyperref}
\usepackage[margin=0.75in]{geometry}
\usepackage{natbib}
\usepackage{xcolor}
\usepackage{algorithm}
\usepackage{algpseudocode}
\usepackage{enumitem}
\usepackage{microtype}
\usepackage{multirow}
\usepackage{subcaption}

\hypersetup{colorlinks=true, linkcolor=blue!60!black, citecolor=blue!60!black, urlcolor=blue!60!black}

\title{Loom: Hybrid Retrieval-Scoring Outfit Recommendation\\with Semantic Material Compatibility\\and Occasion-Aware Embedding Priors}

\author{
  Anushree Berlia\\
  \texttt{github.com/anushreeberlia/loom}
}

\date{May 2026}

\begin{document}
\maketitle

% ══════════════════════════════════════════════════════════════════════════════
\begin{abstract}
We present \textbf{Loom}, an outfit recommendation system that combines
neural embedding retrieval with structured domain scoring to generate
complete, coherent outfits from fashion catalogs.  Given an anchor clothing
item, Loom retrieves complementary pieces via slot-constrained approximate
nearest neighbor search over FashionCLIP embeddings, then scores candidate
outfits using a multi-objective function that integrates six signals:
embedding similarity, color harmony, formality consistency, occasion
coherence, style direction, and within-outfit diversity.

We introduce three techniques that address limitations of purely learned
or purely rule-based approaches:
(1)~\emph{semantic material weight}, which uses CLIP embedding geometry to
infer garment heaviness for layer compatibility without hand-coded material
taxonomies;
(2)~\emph{vibe/anti-vibe occasion priors}, which embed prose descriptions of
occasion contexts as anchor vectors in CLIP space and score items by
differential affinity; and
(3)~\emph{tag-level mood scoring}, which matches individual style tags
(rather than full item descriptions) against mood queries, exploiting the
finding that per-tag embeddings separate moods more cleanly than averaged
item-level embeddings.

The system operates in two deployment modes: a personal wardrobe app with
weather-adaptive, mood-responsive daily suggestions and a learned taste
model, and a Shopify storefront integration with incremental catalog updates.
Full outfit generation completes in under 3~seconds, with vector retrieval
at $\sim$5\,ms over 500-item catalogs.

\smallskip\noindent
Code: \url{https://github.com/anushreeberlia/loom}
\end{abstract}

% ══════════════════════════════════════════════════════════════════════════════
\section{Introduction}
\label{sec:intro}

Outfit recommendation is a compositional task: the system must assemble a
\emph{set} of garments that are mutually compatible across visual style,
color palette, formality level, and contextual appropriateness.  This
distinguishes it from single-item recommendation, where the goal is to find
items similar (or complementary) to a query, and from fashion retrieval,
which surfaces items matching a description.

Prior work has approached outfit compatibility through learned embeddings
\citep{vasileva2018learning, chia2022fashionclip}, graph neural networks on
co-purchase data \citep{cui2019dressing}, and transformer-based set
completion \citep{sarkar2023outfittransformer}.  These methods capture
statistical co-occurrence patterns but struggle with hard constraints that
fashion practitioners consider non-negotiable: a formal blazer should not
appear with gym sneakers regardless of embedding proximity; a silk blouse
should not be layered under a cotton t-shirt; a ``work'' outfit should not
include a sequined clutch.

We present \textbf{Loom}, a hybrid system that uses neural embeddings for
retrieval and combines them with structured domain signals for scoring.  Our
contributions:

\begin{enumerate}[leftmargin=*,topsep=2pt,itemsep=1pt]
\item \textbf{Hybrid retrieval-scoring pipeline}: pgvector ANN retrieval
  with slot constraints, followed by multi-objective combinatorial scoring
  that bridges neural similarity and rule-based fashion domain knowledge.

\item \textbf{Semantic material weight}: We compute garment ``heaviness''
  by measuring cosine distance between an item's text embedding and
  pre-embedded context vectors for ``heavy'' and ``light'' materials.  This
  enables layer compatibility checking without hand-coded material
  taxonomies (Section~\ref{sec:material}).

\item \textbf{Vibe/anti-vibe occasion priors}: We embed prose descriptions
  of occasion ``vibes'' (e.g., \emph{``professional office business
  conservative\ldots''}) and ``anti-vibes'' (e.g., \emph{``sexy revealing
  clubbing\ldots''}) as anchor vectors in CLIP space, scoring items by
  differential affinity with occasion-specific penalty shaping
  (Section~\ref{sec:occasion}).

\item \textbf{Tag-level mood scoring}: We show that matching individual
  style tags against mood queries outperforms matching full item
  descriptions, because CLIP averages away discriminative tag signal in
  longer texts (Section~\ref{sec:tag_mood}).

\item \textbf{Blended multimodal embeddings}: $0.7\;\text{image} +
  0.3\;\text{text}$ metadata fusion in a single 512-d vector, where the
  text component resolves visual ambiguities (e.g., athletic vs.\ casual top
  on a hanger) (Section~\ref{sec:embedding}).
\end{enumerate}

% ══════════════════════════════════════════════════════════════════════════════
\section{Related Work}
\label{sec:related}

\paragraph{Outfit compatibility learning.}
Type-aware embeddings \citep{vasileva2018learning} project items into
type-conditioned spaces.  \citet{cui2019dressing} use graph neural networks
over fashion items.  OutfitTransformer \citep{sarkar2023outfittransformer}
applies self-attention over item sets.  DiFashion \citep{xu2024difashion}
uses diffusion models for outfit generation.  CSA-Net \citep{lin2020csanet}
applies category-based subspace attention.  Our system differs by combining
a learned embedding space (FashionCLIP) with explicit rule-based
constraints---capturing domain knowledge that is difficult to learn from
co-occurrence alone.

\paragraph{Fashion embeddings.}
FashionCLIP \citep{chia2022fashionclip} adapts CLIP \citep{radford2021clip}
to fashion image--text pairs, producing embeddings that capture stylistic
affinity rather than visual similarity.  We use FashionCLIP as our embedding
backbone and build all semantic scoring mechanisms (material weight, occasion
priors, mood matching) on its vector space geometry.

\paragraph{Context-aware recommendation.}
Weather-aware \citep{chen2019weather}, occasion-aware
\citep{lin2020occasion}, and body-shape-aware \citep{hidayati2021dress}
fashion recommendation has been explored.  Our system integrates weather,
mood, and occasion signals through a unified embedding-based scoring
framework rather than separate specialized models.

\paragraph{Production systems.}
Industrial outfit engines at Stitch Fix \citep{chen2019stitchfix}, Amazon
\citep{kang2017visually}, and Zalando operate at scale on large catalogs.
Loom targets small-to-medium catalogs (hundreds to low thousands) where
exhaustive combinatorial scoring is feasible, and unifies personal wardrobe
and B2B retail under the same engine.

% ══════════════════════════════════════════════════════════════════════════════
\section{System Architecture}
\label{sec:architecture}

\begin{figure*}[t]
\centering
\fbox{\parbox{0.95\textwidth}{\centering
\textbf{Loom Pipeline}\\[6pt]
\texttt{Anchor Image} $\;\xrightarrow{\text{Fashion Florence}}\;$
\texttt{Structured Tags} $\;\xrightarrow{\text{FashionCLIP}}\;$
\texttt{512-d Embedding}\\[4pt]
$\downarrow$\\[4pt]
$\forall\, s \in \{\text{bottom, shoes, layer, accessory}\}:$\quad
\texttt{Query}$_s$ $\;\xrightarrow{\text{pgvector HNSW}}\;$
\texttt{Candidates}$_s$ $\;\xrightarrow{\text{Occasion filter}}\;$
\texttt{Filtered}$_s$\\[4pt]
$\downarrow$\\[4pt]
\texttt{Combinatorial Candidates} $\;\xrightarrow{\text{Multi-signal score}}\;$
\texttt{Best outfit per direction}\\[4pt]
$\downarrow$\\[4pt]
\texttt{3 Outfits:}\quad Classic $\mid$ Trendy $\mid$ Bold
}}
\caption{End-to-end pipeline.  For each of three style directions, the
system constructs slot-specific queries, retrieves candidates via cosine ANN
search, filters by occasion, generates combinatorial outfit candidates, and
selects the highest-scoring outfit per direction.}
\label{fig:pipeline}
\end{figure*}

Loom processes an anchor item through four stages:

\begin{enumerate}[leftmargin=*,topsep=2pt,itemsep=1pt]
\item \textbf{Vision}: Fashion Florence \citep{berlia2026florence} extracts
  structured attributes (category, color, material, style, occasion) from the
  item image.
\item \textbf{Embedding}: FashionCLIP produces a 512-d blended
  image--text embedding (Section~\ref{sec:embedding}).
\item \textbf{Retrieval}: Per-slot cosine ANN queries with category
  constraints, occasion filtering, and color-aware reranking
  (Section~\ref{sec:retrieval}).
\item \textbf{Scoring}: Combinatorial candidates are evaluated by
  a composite scoring function and the best outfit per direction is returned
  (Section~\ref{sec:scoring}).
\end{enumerate}

% ══════════════════════════════════════════════════════════════════════════════
\section{Method}
\label{sec:method}

\subsection{Multimodal Item Representation}
\label{sec:embedding}

Each item is embedded into a 512-dimensional vector space using FashionCLIP
\citep{chia2022fashionclip}, run as ONNX on CPU.  The final embedding
blends image and text signals:

\begin{equation}
  \mathbf{e} = \text{normalize}\big(
    \alpha \cdot \mathbf{e}_{\text{img}} + \beta \cdot \mathbf{e}_{\text{txt}}
  \big)
\label{eq:blend}
\end{equation}

\noindent where $\alpha = 0.7$, $\beta = 0.3$.  The text component
$\mathbf{e}_{\text{txt}}$ encodes a natural-language description constructed
from the item's attributes (e.g., \texttt{"gray sporty fitted polyester
workout top"}).

\paragraph{Motivation.}
Image-only embeddings cannot distinguish an athletic top from a casual top
when both look similar on a hanger.  The text component provides semantic
identity that resolves such visual ambiguities.  The 70/30 weighting
preserves visual fidelity while injecting attribute-level discrimination.

\subsection{Slot-Constrained Retrieval with Occasion Priors}
\label{sec:retrieval}

For each outfit slot $s \in \{\text{top, bottom, shoes, layer, accessory}\}$,
we construct a slot-specific query embedding $\mathbf{q}_s$ from the anchor
item's attributes, slot-specific item type hints, and directional style cues.

\paragraph{Query construction.}
Each query concatenates: (i)~the anchor item's color and style tags,
(ii)~a slot hint (e.g., \texttt{"women's skirt, jeans, trousers"}), (iii)~a
directional modifier (e.g., \texttt{"classic timeless polished"}).  Queries
are batch-embedded via the FashionCLIP text encoder.

\paragraph{ANN retrieval.}
Embeddings are stored in PostgreSQL with pgvector using HNSW indexing
($m{=}16$, $\texttt{ef\_construction}{=}64$, cosine distance).  Each slot
query retrieves $2k$ candidates filtered by:
\begin{equation}
  \mathcal{C}_s = \text{top-}2k\big(
    \{\mathbf{e}_i : c_i = s, \; i \notin \mathcal{E}\},\;
    d_{\cos}(\mathbf{q}_s, \mathbf{e}_i)
  \big)
\end{equation}
where $c_i$ is item category and $\mathcal{E}$ excludes the anchor and
previously selected items.

\paragraph{Soft color reranking.}
Retrieved candidates are reranked with distance multipliers:
preferred colors receive $0.85\times$ distance, neutrals $0.9\times$.
Multiplicative noise $\epsilon \sim \mathcal{U}(0.95, 1.05)$ is applied to
distances to promote diversity across repeated calls.

\subsubsection{Occasion-Aware Filtering}
\label{sec:occasion}

Rather than mapping occasions to discrete categories, we embed rich prose
descriptions as anchor vectors in CLIP space.

\paragraph{Vibe/anti-vibe vectors.}
For each occasion $o$, we pre-embed two context strings:
\begin{align}
  \mathbf{v}_o &= \text{embed}(\text{vibe}_o) \\
  \mathbf{a}_o &= \text{embed}(\text{anti-vibe}_o)
\end{align}

For example, the ``work'' occasion uses:
\begin{quote}
\small
\textbf{vibe}: \textit{``professional office business conservative modest
polished refined tailored structured classic understated\ldots''}\\
\textbf{anti-vibe}: \textit{``sexy revealing provocative clubbing nightlife
party halter tank top cami strappy\ldots mini skirt crop top\ldots''}
\end{quote}

\paragraph{Scoring.}
Each item is scored by differential affinity with occasion-specific penalty
scaling:
\begin{equation}
  S_{\text{occ}}(i, o) = \cos(\mathbf{e}_i, \mathbf{v}_o) -
    \lambda_o \cdot \max\!\big(0,\;
    \cos(\mathbf{e}_i, \mathbf{a}_o) - \cos(\mathbf{e}_i, \mathbf{v}_o)\big)
\label{eq:occasion}
\end{equation}
where $\lambda_o$ controls strictness: $\lambda_{\text{work}} = 3.0$ (strict
filtering), $\lambda_{\text{going-out}} = \lambda_{\text{smart-casual}} =
3.0$, $\lambda_{\text{casual}} = 1.0$.  For work, an additional strong
penalty $2.0 \cdot \cos(\mathbf{e}_i, \mathbf{a}_o)$ is applied
unconditionally.

\paragraph{Threshold filtering.}
Candidates are kept if they score within a fraction of the top score:
85\% for work (strict), 70\% for casual (permissive), with a minimum floor
of 3 candidates.

\subsection{Semantic Material Compatibility}
\label{sec:material}

Layering requires that outer garments be heavier than inner ones: a blazer
over a silk blouse is correct; a cotton t-shirt over a wool coat is not.
Rather than maintaining a hand-coded material weight table, we compute
material weight semantically.

\paragraph{Context vectors.}
Two context strings are pre-embedded:
\begin{align}
  \mathbf{h} &= \text{embed}(\text{``thick heavy warm winter chunky\ldots''}) \\
  \mathbf{l} &= \text{embed}(\text{``thin light airy breathable summer\ldots''})
\end{align}

\paragraph{Weight score.}
For each item $i$, we construct a text string from its material, category,
and weight-relevant name keywords, embed it, and compute:
\begin{equation}
  w_i = \cos(\mathbf{e}_{\text{mat}(i)}, \mathbf{h}) -
    \cos(\mathbf{e}_{\text{mat}(i)}, \mathbf{l})
  \;\in [-1, +1]
\label{eq:material_weight}
\end{equation}

\paragraph{Compatibility check.}
A layer is compatible with a top if $w_{\text{layer}} \geq
w_{\text{top}} - 0.15$, with the tolerance absorbing embedding noise.

\subsection{Multi-Objective Outfit Scoring}
\label{sec:scoring}

Given an anchor item and candidate pieces per slot, outfits are scored by a
composite function.

\subsubsection{Intent Vector}

The intent vector combines the anchor embedding with optional taste signals:
\begin{equation}
  \mathbf{v} = \mathbf{e}_{\text{anchor}} + \gamma \cdot
    \mathbf{t}_{\text{like}} - \delta \cdot \mathbf{t}_{\text{dislike}}
\label{eq:intent}
\end{equation}
where $\gamma = 0.15$, $\delta = 0.05$, and
$\mathbf{t}_{\text{like}}, \mathbf{t}_{\text{dislike}} \in \mathbb{R}^{512}$
are exponential moving averages of liked/disliked outfit embeddings.

\subsubsection{Weighted Similarity}

Each piece is scored against the intent vector with slot-dependent weights:
\begin{equation}
  S_{\text{sim}} = \sum_{s \in \mathcal{S}} w_s \cdot
    \cos(\mathbf{v}, \mathbf{e}_s) + w_{\bot\text{-shoe}} \cdot
    \cos(\mathbf{e}_{\text{bot}}, \mathbf{e}_{\text{shoe}})
\label{eq:sim}
\end{equation}
with $w_{\text{bottom}} = 0.35$, $w_{\text{shoes}} = 0.25$,
$w_{\text{acc}} = 0.15$, $w_{\text{layer}} = 0.10$,
$w_{\text{top}} = 0.10$, $w_{\bot\text{-shoe}} = 0.05$.

The asymmetric weighting reflects the empirical importance of bottom and shoe
choices to overall outfit coherence: a wrong bottom or wrong shoes breaks
an outfit more than a suboptimal accessory.

\subsubsection{Style Direction Bonuses}

Three scoring profiles produce stylistic diversity:

\begin{table}[h]
\centering
\small
\caption{Style direction configurations.}
\label{tab:directions}
\begin{tabular}{llll}
\toprule
Direction & Style tags & Color policy & Vibe \\
\midrule
Classic & classic, minimalist, & neutrals & timeless, \\
        & elegant, preppy & & polished \\
Trendy  & streetwear, chic, & two-tone & modern, \\
        & statement, casual & & forward \\
Bold    & edgy, statement, & contrast & daring, \\
        & romantic, bohemian & & eye-catching \\
\bottomrule
\end{tabular}
\end{table}

Each direction contributes a bonus $B_{\text{dir}} \in [0, 0.3]$ based on
tag matches, color policy adherence, and direction-specific rules (e.g.,
Classic rewards neutral shoes/accessories; Bold rewards brightness contrast
with the anchor; Trendy penalizes same-neutral bottom--shoes).

\subsubsection{Constraint Penalties}

\paragraph{Color harmony.}
$P_{\text{color}} = 0.1 \cdot \max(0, n_{\text{non-neutral}} - 2)$ penalizes
outfits with too many competing colors.  Hard violations (repeated
non-neutral color from anchor) immediately set score to $-1$.

\paragraph{Formality consistency.}
Items are assigned formality levels $\ell \in \{0, 1, 2\}$ (casual,
smart-casual, dressy) via keyword matching.  Penalty:
$P_{\text{form}} = 0.2 \cdot \max(0, \Delta\ell - 1)$.

\paragraph{Occasion coherence.}
$P_{\text{occ}} = 0.15$ if the intersection of all items' inferred occasion
tag sets is empty.

\paragraph{Diversity control.}
Multiple statement pieces: $P_{\text{div}} = 0.1 \cdot
(n_{\text{statement}} - 1)$.  Color family analysis (warm/cool/neutral)
contributes harmony bonuses $B_{\text{harm}} \in [-0.05, +0.05]$.

\subsubsection{Total Score}

\begin{equation}
  S = S_{\text{sim}} + B_{\text{dir}} + B_{\text{harm}}
    - P_{\text{color}} - P_{\text{form}} - P_{\text{occ}} - P_{\text{div}}
\label{eq:total}
\end{equation}

\subsubsection{Candidate Generation}

For each direction, the top $k{=}3$ candidates per slot yield combinatorial
outfit candidates (capped at 8 via sampling from
$\prod_{s} |\mathcal{C}_s|$).  The highest-scoring candidate per direction
is selected.

\subsection{Tag-Level Mood Scoring}
\label{sec:tag_mood}

When users provide free-text mood queries (e.g., ``job interview,'' ``beach
day''), the system must rank items by relevance to the mood.

\paragraph{Observation.}
Embedding full item descriptions (``gray sporty fitted polyester top'')
against mood queries produces nearly uniform similarity scores across items,
because CLIP averages away discriminative signal in longer texts.  However,
individual style tags show clear separation:

\vspace{4pt}
\begin{center}
\small
\begin{tabular}{lcc}
\toprule
Tag & $\cos(\text{tag}, \text{``workout''})$ \\
\midrule
``sporty'' & 0.69 \\
``casual'' & 0.64 \\
``classic'' & 0.58 \\
\bottomrule
\end{tabular}
\end{center}
\vspace{4pt}

\paragraph{Method.}
For each item, we compute per-tag similarity to the mood and take the
maximum:
\begin{equation}
  S_{\text{mood}}(i) = \max_{t \in \text{style\_tags}(i)}
    \cos\!\big(\text{embed}(t),\; \text{embed}(\text{mood})\big)
\end{equation}

This uses only \texttt{style\_tags} (not \texttt{occasion\_tags}), because
generic occasion tags like ``everyday'' inflate casual scores uniformly.

% ══════════════════════════════════════════════════════════════════════════════
\section{Personalization}
\label{sec:personal}

\subsection{Taste Vectors}

User feedback (like/dislike) updates taste vectors via exponential moving
average:
\begin{equation}
  \mathbf{t}^{(n+1)} = \frac{n}{n+1} \cdot \mathbf{t}^{(n)}
    + \frac{1}{n+1} \cdot \bar{\mathbf{e}}_{\text{outfit}}
\end{equation}
where $\bar{\mathbf{e}}_{\text{outfit}}$ is the mean embedding of outfit
items.  Taste vectors are blended into retrieval queries ($\gamma_{\text{ret}}
= 0.25$, $\delta_{\text{ret}} = 0.15$) and into scoring via the intent
vector (Eq.~\ref{eq:intent}).

\subsection{Weather Adaptation}

The system queries OpenWeatherMap and adjusts generation: rain avoids suede,
cold forces the layer slot and boosts heavy materials (via semantic weight,
Eq.~\ref{eq:material_weight}), heat skips layers and prefers light fabrics.

\subsection{Rotation Queue}

A per-user, per-occasion FIFO queue tracks when each item was last
suggested.  Recency penalties: $-20$ (today), $-10$ (within 3 days), $-5$
(within 7 days).  Users can mark outfits as ``worn,'' further deprioritizing
those items.

% ══════════════════════════════════════════════════════════════════════════════
\section{Shopify Deployment}
\label{sec:shopify}

\subsection{Incremental Catalog Updates}

When merchants add products, Shopify webhooks trigger selective
reprocessing.  A category-based invalidation map determines which cached
outfits must be regenerated:

\begin{table}[h]
\centering
\small
\caption{Invalidation map: adding category $c$ invalidates cached outfits
for anchor categories $\mathcal{A}(c)$.}
\label{tab:invalidation}
\begin{tabular}{ll}
\toprule
New item $c$ & Invalidated anchors $\mathcal{A}(c)$ \\
\midrule
top       & bottom, shoes, dress \\
bottom    & top, shoes \\
shoes     & top, bottom, dress \\
layer     & top, bottom, shoes, dress \\
accessory & top, bottom, shoes, dress, layer \\
dress     & shoes \\
\bottomrule
\end{tabular}
\end{table}

This avoids full-catalog regeneration: new shoes only invalidate outfits
for anchors where shoes appear as a component slot.

\subsection{Storefront Rendering}

A Shopify theme extension renders ``Shop the Look'' blocks on product pages
with pre-computed outfit collages, item details, prices, and purchase links.

% ══════════════════════════════════════════════════════════════════════════════
\section{Experiments}
\label{sec:experiments}

\subsection{Polyvore Benchmark}

We evaluate on the Polyvore Outfits dataset \citep{vasileva2018learning},
the standard benchmark for outfit compatibility:

\paragraph{Tasks.}
(i)~\textbf{Fill-in-the-blank (FITB)}: given an incomplete outfit and 4
candidates, select the correct missing item.
(ii)~\textbf{Compatibility prediction}: binary classification of whether
a set of items forms a compatible outfit (AUC).

\begin{table}[h]
\centering
\small
\caption{Polyvore Outfits benchmark (prior work).  Loom is not directly
evaluated on FITB/compatibility tasks because its task formulation differs:
complete outfit \emph{generation} from a single anchor, not pairwise
compatibility scoring or item retrieval.}
\label{tab:polyvore}
\begin{tabular}{lcc}
\toprule
Method & FITB Acc.\ (\%) & Compat.\ AUC \\
\midrule
Type-Aware \citep{vasileva2018learning} & 53.7 & 0.86 \\
CSA-Net \citep{lin2020csanet} & 63.7 & 0.91 \\
OutfitTransformer \citep{sarkar2023outfittransformer} & 67.1 & 0.93 \\
DiFashion \citep{xu2024difashion} & -- & -- \\
\bottomrule
\end{tabular}
\end{table}

\subsection{Ablation Study}

We measure the contribution of each component on outfit quality, evaluated
on our internal catalog with both automatic metrics and human evaluation:

\begin{table}[h]
\centering
\small
\caption{Ablation results (50 anchors $\times$ 3 directions = 150 outfits
per configuration, 620-item catalog).}
\label{tab:ablation}
\begin{tabular}{lccc}
\toprule
Configuration & Score & $\Delta$ & Viol.\,\% \\
\midrule
Full system              & 0.218 & ---      &  6.0 \\
$-$ Blended embeddings   & 0.163 & $-$0.055 & 10.0 \\
$-$ Occasion filtering   & 0.165 & $-$0.053 & 10.7 \\
$-$ Material compat.     & 0.186 & $-$0.032 &  8.7 \\
$-$ Distance noise       & 0.145 & $-$0.073 & 12.0 \\
$-$ Direction reranking  & 0.022 & $-$0.196 & 10.0 \\
$-$ Formality check      & 0.205 & $-$0.013 &  8.7 \\
\midrule
Random retrieval         & 0.125 & $-$0.093 & 14.0 \\
\bottomrule
\end{tabular}
\end{table}

\subsection{Human Evaluation}

\textit{[Planned: N=50--100 outfit evaluations by 3--5 annotators, rating
coherence (1--5), occasion-appropriateness (1--5), diversity across 3
directions (1--5), and ``would you wear this?'' (binary).  Full vs.\
ablated variants.]}

\subsection{Diversity Metrics}

We measure diversity across the three style directions per anchor item:

\begin{itemize}[leftmargin=*,topsep=2pt,itemsep=1pt]
\item \textbf{Color spread}: number of distinct color families across the
  3 outfits.
\item \textbf{Subtype diversity}: proportion of distinct item subtypes
  (e.g., ``heels'' vs.\ ``sneakers'' for shoes) across directions.
\item \textbf{Embedding variance}: mean pairwise cosine distance between
  corresponding slots across the 3 directions.
\end{itemize}

\subsection{Efficiency}

\begin{table}[h]
\centering
\small
\caption{Latency breakdown (500-item catalog, CPU inference).}
\label{tab:latency}
\begin{tabular}{lr}
\toprule
Stage & Latency \\
\midrule
Vision (Fashion Florence, remote) & $\sim$1{,}500 ms \\
Embedding (FashionCLIP ONNX, local) & $\sim$200 ms \\
Retrieval (pgvector HNSW, 4 slots) & $\sim$5 ms \\
Scoring (3 directions $\times$ 8 candidates) & $\sim$50 ms \\
\midrule
\textbf{Total} & $<$\textbf{3{,}000 ms} \\
\bottomrule
\end{tabular}
\end{table}

% ══════════════════════════════════════════════════════════════════════════════
\section{Discussion}
\label{sec:discussion}

\paragraph{Similarity vs.\ complementarity.}
FashionCLIP embeddings capture stylistic \emph{similarity}: items close in
vector space share visual and semantic characteristics.  Outfit composition,
however, requires \emph{complementarity}---a navy blazer should retrieve
khaki chinos, not another navy blazer.  Our scoring function approximates
complementarity through slot constraints and color harmony rules, but a
learned complement projection $f: \mathbb{R}^{512} \to \mathbb{R}^{512}$
that maps an item's embedding to its ideal complement would more directly
encode this relationship.

\paragraph{Scalability.}
The combinatorial scoring phase ($O(k^{|\mathcal{S}|})$) is feasible for
small-to-medium catalogs but would not scale to millions of items.  For
larger catalogs, hierarchical pruning or learned scoring models would be
needed to replace the exhaustive approach.

\paragraph{Occasion prior coverage.}
The vibe/anti-vibe approach requires prose descriptions for each occasion.
Extending to new occasions (e.g., ``funeral,'' ``job interview at a
startup'') requires only writing new context strings---no retraining---but
the quality of scoring depends on the prose descriptions capturing the right
associations.

% ══════════════════════════════════════════════════════════════════════════════
\section{Future Work}
\label{sec:future}

\begin{enumerate}[leftmargin=*,topsep=2pt,itemsep=1pt]
\item \textbf{Complement embeddings}: Train a projection head
  $\mathbf{e}_{\text{item}} \to \mathbf{e}_{\text{complement}}$ using
  styled outfit co-occurrence data.
\item \textbf{Cross-catalog shopping}: Enable users to discover and purchase
  complementary retail items to fill wardrobe gaps.
\item \textbf{Learned scoring}: Replace or augment the rule-based scoring
  components with a learned outfit scorer trained on human preference data.
\item \textbf{Visual rendering}: Move from flat collage grids to layered
  or mannequin-style composition views.
\end{enumerate}

% ══════════════════════════════════════════════════════════════════════════════
\section{Conclusion}

We presented Loom, a hybrid retrieval-scoring system for outfit
recommendation that combines FashionCLIP embedding retrieval with structured
domain scoring.  Three technical contributions---semantic material weight,
vibe/anti-vibe occasion priors, and tag-level mood scoring---address
specific limitations of purely neural approaches by exploiting CLIP's vector
space geometry for tasks beyond simple similarity.  The system generates
three stylistically distinct outfits in under 3 seconds and is deployed in
both personal wardrobe and Shopify retail contexts.

% ══════════════════════════════════════════════════════════════════════════════
\bibliographystyle{plainnat}
\begin{thebibliography}{20}

\bibitem[Berlia(2026)]{berlia2026florence}
Berlia, A.
\newblock Fashion {Florence}: Fine-tuning {Florence-2} for structured fashion
attribute extraction.
\newblock \emph{arXiv preprint}, 2026.

\bibitem[Chen et~al.(2019a)]{chen2019weather}
Chen, W., Huang, P., Xu, J., et~al.
\newblock Weather-aware fashion recommendation.
\newblock In \emph{WSDM}, 2019.

\bibitem[Chen et~al.(2019b)]{chen2019stitchfix}
Chen, J., Wang, X., et~al.
\newblock Algorithms and data-driven approaches at {Stitch Fix}.
\newblock In \emph{KDD Industry Track}, 2019.

\bibitem[Chia et~al.(2022)]{chia2022fashionclip}
Chia, P.~J., Attanasio, G., Bianchi, F., et~al.
\newblock {FashionCLIP}: Connecting language and images for product
representations.
\newblock \emph{arXiv preprint arXiv:2204.03972}, 2022.

\bibitem[Cui et~al.(2019)]{cui2019dressing}
Cui, Z., Li, Z., Wu, S., Zhang, X., and Wang, L.
\newblock Dressing as a whole: Outfit compatibility learning based on node-wise
graph neural networks.
\newblock In \emph{WWW}, 2019.

\bibitem[Hidayati et~al.(2021)]{hidayati2021dress}
Hidayati, S.~C., Hsu, C.-C., Chang, Y.-T., et~al.
\newblock What dress fits me best? Fashion recommendation on the clothing
fit.
\newblock In \emph{ACM Multimedia}, 2021.

\bibitem[Kang et~al.(2017)]{kang2017visually}
Kang, W.-C., Fang, C., Wang, Z., and McAuley, J.
\newblock Visually-aware fashion recommendation and design with generative
image models.
\newblock In \emph{ICDM}, 2017.

\bibitem[Lin et~al.(2020a)]{lin2020csanet}
Lin, Y., Ren, P., Chen, Z., Ren, Z., Ma, J., and de~Rijke, M.
\newblock Explainable outfit recommendation with joint outfit matching and
comment generation.
\newblock \emph{IEEE TKDE}, 2020.

\bibitem[Lin et~al.(2020b)]{lin2020occasion}
Lin, Y., Moosaei, M., and Yang, H.
\newblock {OutfitNet}: Fashion outfit recommendation with attention-based
multiple instance learning.
\newblock In \emph{WWW}, 2020.

\bibitem[Radford et~al.(2021)]{radford2021clip}
Radford, A., Kim, J.~W., Hallacy, C., et~al.
\newblock Learning transferable visual models from natural language
supervision.
\newblock In \emph{ICML}, 2021.

\bibitem[Sarkar et~al.(2023)]{sarkar2023outfittransformer}
Sarkar, R., Bodla, N., Vasileva, M., et~al.
\newblock {OutfitTransformer}: Learning outfit representations for fashion
recommendation.
\newblock In \emph{WACV}, 2023.

\bibitem[Vasileva et~al.(2018)]{vasileva2018learning}
Vasileva, M., Plummer, B.~A., Dusber, K., et~al.
\newblock Learning type-aware embeddings for fashion compatibility.
\newblock In \emph{ECCV}, 2018.

\bibitem[Xu et~al.(2024)]{xu2024difashion}
Xu, Y., Yan, Y., and Lin, D.
\newblock {DiFashion}: Generative fashion outfit recommendation with
diffusion models.
\newblock In \emph{SIGIR}, 2024.

\end{thebibliography}

\end{document}
